\uuid{eFF8}
\titre{Changements de variables}
\auteur{hanani ; mzari}
\organisation{exo7}
\licence{CC BY-NC-SA 4.0}
\datecreate{2020-01-20}
\niveau{L1}
\module{Analyse}
\chapitre{Calcul d'intégrales}
\sousChapitre{Autre}
\difficulte{3}
\theme{Calculs d'intégrales, changements de variables}
\exo7qcmid{lille-2.integrales.027}

\contenu{
\begin{question}
\qclassification{127.99}
\qtags{motcle=changements de variables}
Parmi les affirmations suivantes, cocher celles qui sont vraies :
\begin{answers}
\good{Le changement de variable $t=\sqrt{x+1}$ donne $\displaystyle \int _{3}^{8}\frac{\mathrm{d}\, x}{x\sqrt{x+1}}=\int _{2}^{3}\frac{2\mathrm{d}\, t}{t^2-1}$.}
\good{$\displaystyle \int _{3}^{8}\frac{\mathrm{d}\, x}{x\sqrt{x+1}}=\ln \frac{3}{2}$.}
\bad{Le changement de variable $t=\sqrt{x+1}$ donne $\displaystyle \int _3^8\frac{\sqrt{x+1}}{x}\mathrm{d}x=\int _2^3\frac{t^2\, \mathrm{d}t}{t^2-1}$.}
\bad{$\displaystyle \int _3^8\frac{\sqrt{x+1}}{x}\mathrm{d}x=1+\frac{1}{2}\ln \frac{3}{2}$.}
\end{answers}
\vskip3mm
\begin{explanations}
Avec $t=\sqrt{x+1}\Rightarrow x=t^2-1$, on a : $\mathrm{d}x=2t\, \mathrm{d}t$, on a : $t(3)=2$, $t(8)=3$ et
$$\displaystyle \int _3^8\frac{\mathrm{d}\, x}{x\sqrt{x+1}}=\int _{2}^{3}\frac{2\mathrm{d}\, t}{t^2-1}=\left[\ln \frac{t-1}{t+1}\right]_{2}^{3}=\ln \frac{3}{2}.$$
De même, $\displaystyle \int _3^8\frac{\sqrt{x+1}}{x}\mathrm{d}x=\int _{2}^{3}\frac{2t^2\, \mathrm{d}t}{t^2-1}$. Or $\displaystyle \frac{2t^2}{t^2-1}=2+\frac{1}{t-1}-\frac{1}{t+1}$, donc
$$\displaystyle \int _3^8\frac{\sqrt{x+1}}{x}\mathrm{d}x=\left[2t+\ln \frac{t-1}{t+1}\right]_{2}^{3}=2+\ln \frac{3}{2}.$$
\end{explanations}
\end{question}
}
